% Part: first-order-logic
% Chapter: lindstrom
% Section: ls-property

\documentclass[../../../include/open-logic-section]{subfiles}

\begin{document}

\olfileid{mod}{lin}{lsp}

\olsection{د متناهي شاهد او لوېنهايم--سکولم خاصيتونه}

اوس د متناهي شاهد او لوېنهايم--سکولم خاصيتونه د انتزاعي منطقونو حالت
ته په څرګند ډول غځوو۔

\begin{defn}
يو انتزاعي منطق~$\tuple{L, \models_L}$ \emph{د متناهي شاهد خاصيت}
لري که د~$L(\Lang{L})$-!!{sentence}s هر سټ~$\Gamma$ هغه وخت د صدق
وړ وي چې هر متناهي~$\Gamma_0 \subseteq \Gamma$ ئې د صدق وړ وي۔
\end{defn}

\begin{defn}
$\tuple{L, \models_L}$ \emph{د لوېنهايم--سکولم ښکته خاصيت} لري که
هر د صدق وړ~$\Gamma$ يو !!a{enumerable} مدل ولري۔
\end{defn}


د~\olref[bas][pis]{defn:partialisom} د جزوي آيزومورفيزم مفهوم سوچه
«الجبري» دے؛ يعنې د ژبې د !!{sentence}s له يادونې پرته او يوازې د
!!{structure}s د~!!{language}~$\Lang{L}$ د ثابتونو له مخې ورکړل شوے
دے۔ نو دا مفهوم پر انتزاعي منطقونو هم عملي کېږي۔ د لومړۍ درجې منطق
په حالت کښې له~\olref[bas][pis]{thm:p-isom2} څخه پوهېږو چې که دوه
!!{structure}s په جزوي ډول آيزومورف وي، نو ابتدايي هم‌ارز دي۔ هغه
ثبوت انتزاعي منطقونو ته نۀ غځېږي، ځکه پر !!{formula}s استقرا د خپل
سري~$!E \in L(\Lang{L})$ دپاره ښايي موجوده نۀ وي؛ خو که د
لوېنهايم--سکولم خاصيت صدق وکړي، قضيه بيا هم رښتونې ده۔

\begin{thm}
\ollabel{thm:abstract-p-isom}
فرض کړئ~$\tuple{L, \models_L}$ يو نورمال منطق دے چې د
لوېنهايم--سکولم خاصيت لري۔ بيا هر دوه په جزوي ډول آيزومورف
!!{structure}s په~$\tuple{L, \models_L}$ کښې ابتدايي هم‌ارز دي۔
\end{thm}

\begin{proof}
فرض کړئ~$\Struct{M} \simeq_p \Struct{N}$، خو د کوم~$!E$ دپاره
$\Struct{M} \models_L !E$ او~$\Struct{N} \not\models_L !E$ هم وي۔ د
آيزومورفيزم د خاصيت له مخې فرض کولے شو چې~$\Domain{M}$ او
$\Domain{N}$ سره بېل دي؛ او د غځونې د خاصيت له مخې فرض کولے شو چې
$!E \in L(\Lang{L})$ د يوې متناهي !!{language}~$\Lang{L}$ دپاره
وي۔ $\mathcal{I}$ دې د~$\Struct{M}$ او~$\Struct{N}$ تر منځ د جزوي
آيزومورفيزمونو يو سټ وي؛ له عموميت څخه بې له زيان څخه دا هم فرض کړئ
چې که~$p \in \mathcal{I}$ او~$q \subseteq p$ وي، نو
$q \in \mathcal{I}$ هم وي۔

$\Domain{M}^{<\omega}$ د~$\Domain{M}$ د !!{element}s د متناهي لړيو
سټ دے۔ $S$ دې پر~$\Domain{M}^{<\omega}$ هغه درې‌ځاييزه اړيکه وي چې
نښلون ښيي؛ يعنې که~$\mathbf{a}, \mathbf{b},
\mathbf{c} \in \Domain{M}^{<\omega}$ وي، نو
$S(\mathbf{a}, \mathbf{b}, \mathbf{c})$ هله او يوازې هله صدق کوي
چې~$\mathbf{c}$ د~$\mathbf{a}$ او~$\mathbf{b}$ نښلون وي۔ همدارنګه،
$T$ دې هغه درې‌ځاييزه اړيکه وي چې د~$b \in M$ او
$\mathbf{a}, \mathbf{c} \in \Domain{M}^{<\omega}$ دپاره
$T(\mathbf{a}, b, \mathbf{c})$ هله او يوازې هله صدق کوي چې
$\mathbf{a} = a_1, \dots a_n$ او
$\mathbf{c} = a_1, \dots a_n, b$ وي۔ دوه نوي درې‌ځاييز
!!{predicate}s~$P$ او~$Q$ واخلئ، او !!{structure}~$\Struct{M}^*$
داسې جوړ کړئ چې جهان ئې~$\Domain{M} \cup \Domain{M}^{<\omega}$ وي،
$\Struct{M}$ د فرعي جوړښت په توګه ولري، او~$P$ او~$Q$ په نښلوونکو
اړيکو~$S$ او~$T$ تفسير کړي؛ نو~$\Struct{M}^*$ د
!!{language}~$\Lang{L} \cup \{ P, Q\}$ جوړښت دے۔ د اړتيا پر وخت د
جهان د دوو برخو بېل نقلونه اخلو، څو دا اتحاد سره بېل کوډ شوے وي۔

$\Domain{N}^{<\omega}$،~$S'$،~$T'$،~$P'$،~$Q'$ او
$\Struct{N}^*$ په ورته ډول تعريف کړئ۔ ځکه چې د فرض له مخې
$\Struct{M} \simeq_p \Struct{N}$، د~$\Domain{M}^{<\omega}$ او
$\Domain{N}^{<\omega}$ تر منځ داسې اړيکه~$I$ شته چې
$I(\mathbf{a}, \mathbf{b})$ هله او يوازې هله صدق کوي چې
$\mathbf{a}$ او~$\mathbf{b}$ آيزومورف وي او د
\olref[bas][pis]{defn:partialisom} د مخ او شا شرط پوره کړي۔ اوس
$\Struct{K}$ هغه !!{structure} ونيسئ چې !!{domain} ئې د
$\Struct{M}^*$ او~$\Struct{N}^*$ د !!{domain}s سره بېل کوډ شوے اتحاد
وي،~$\Struct{M}^*$ او~$\Struct{N}^*$ د فرعي !!{structure}s په توګه
ولري، او !!{language} ئې يو اضافي دوه‌ځاييز !!{predicate}~$R$ ولري چې په
اړيکه~$I$ تفسيرېږي، او داسې !!{predicate}s هم ولري چې د دواړو برخو
!!{domain}s~$\Domain{M}^*$ او~$\Domain{N}^*$ ټاکي۔

\begin{figure}[h]
  \centering
  \begin{tikzpicture}[node distance=2cm, auto, thick, >=stealth']
    \draw [rounded corners] (0,0) -- (8,0) -- (8,4) -- (0,4) --  cycle;
    \draw (2,2) circle (0.5cm);
    \draw (2,2) circle (1.25cm);
    \draw (6,2) circle (0.5cm);
    \draw (6,2) circle (1.25cm);
    \path node at (0.75,3.5) {\large $\Struct{K}$};
    \path node at (2,2) {\large $\Struct{M}$};
    \path node at (6,2) {\large $\Struct{N}$};
    \path node at (3.5,1) {\large $\Struct{M}^*$};
    \path node at (7.5,1) {\large $\Struct{N}^*$};
    \node (Idom) at (2.8,2) {};
    \node (Irng) at (5.2,2) {};
    \draw[<->, bend left] (Idom) to node {\large $I$} (Irng) ;
  \end{tikzpicture}
  \caption{!!{structure}~$\Struct{K}$ او د هغه دننه جزوي
    آيزومورفيزم۔}
\end{figure}

مهمه کتنه دا ده چې د~!!{structure}~$\Struct{K}$ په !!{language} کښې يوه د
انتزاعي منطق !!{sentence}~$!D_1$ شته چې په~$\Struct{K}$ کښې رښتونې
ده او وايي~$\Struct{M} \models_L !E$ او
$\Struct{N} \not\models_L !E$؛ د دې دپاره د نسبي کولو خاصيت پکار
دے۔ همدارنګه، يوه د لومړۍ درجې !!{sentence}~$!D_2$ په~$\Struct{K}$
کښې رښتونې ده او وايي چې~$\Struct{M} \simeq_p \Struct{N}$ د جزوي
آيزومورفيزم~$I$ له لارې۔ د لوېنهايم--سکولم د خاصيت له مخې،~$!D_1$
او~$!D_2$ په ګډه په يوه !!a{enumerable} مدل~$\Struct{K}_0$ کښې
رښتونې دي؛ په دې مدل کښې په جزوي ډول آيزومورف فرعي جوړښتونه
$\Struct{M}_0$ او~$\Struct{N}_0$ شته، داسې چې
$\Struct{M}_0 \models_L !E$ او~$\Struct{N}_0 \not\models_L !E$۔ خو
!!{enumerable} او په جزوي ډول آيزومورف !!{structure}s د
\olref[bas][pis]{thm:p-isom1} له مخې په حقيقت کښې آيزومورف دي؛ دا د
نورمالو انتزاعي منطقونو د آيزومورفيزم له خاصيت سره تناقض دے۔
\textbf{د ژباړې د نسخې سمون (OLMOD-026):} سرچينې دواړه شاهدې جملې د
لومړۍ درجې بللې وې، حال دا چې لومړۍ جمله د خپل سري انتزاعي جملې د
نسبي کولو له لارې لاس ته راځي؛ دلته لومړۍ د انتزاعي منطق جمله او
دويمه د لومړۍ درجې جمله وبلل شوه۔
\textbf{د ژباړې د نسخې سمون (OLMOD-027):} سرچينې ګډ جوړښت ته د
پخواني چپ جوړښت هماغه نوم، او بيا ښکته مدل او چپ فرعي جوړښت ته هم
يو شان نوم ورکړے و؛ دلته ګډ جوړښت او د هغه شمېرېدونکی مدل په بېلو
نښو ښودل شوي، او په انځور کښې بهرنی نوم هم ورسره سم شوے دے۔
\end{proof}

\end{document}
